\documentclass[11pt]{article}
\usepackage[margin=2.1cm]{geometry}
\usepackage{booktabs}
\usepackage{amsmath}
\usepackage{enumitem}
\usepackage[T1]{fontenc}
\usepackage{lmodern}

\setlength{\parindent}{0pt}
\setlength{\parskip}{5pt}
\newcommand{\code}[1]{\texttt{#1}}
\newcommand{\answerbox}[1]{\vspace{2pt}\fbox{\parbox[c][#1][c]{\dimexpr\linewidth-2\fboxsep-2\fboxrule}{\ }}\vspace{4pt}}

\begin{document}

\begin{center}
{\Large\bfseries EE\,362 --- Pre-Lab 7}\\[2pt]
{\large Next Instruction Address Generation I}\\[6pt]
\begin{tabular}{@{}p{6.5cm}p{4cm}p{4cm}@{}}
Name: \hrulefill & Section: Mon / Tue & Date: \hrulefill
\end{tabular}
\end{center}

\vspace{2pt}\hrule\vspace{4pt}

{\small\textbf{Due at the start of the lab session.} Five points, about twenty
minutes, on paper. \textbf{Bring a calculator or do the hex by hand} --- this
one is arithmetic.}

\vspace{4pt}\hrule

\section*{Question 1 --- Paper first}

Here is a real instruction from the program your CPU will run in Lab 12:

\begin{center}
\begin{tabular}{lll}
\toprule
address & encoding & instruction \\
\midrule
\code{0x0000000C} & \code{0x18200002} & \code{BLEZ r1, +2} \\
\bottomrule
\end{tabular}
\end{center}

\textbf{(a)} Compute the branch target. Show all three steps.

\vspace{4pt}
\code{NSA} $= $ \rule{4cm}{0.4pt} \hfill
offset $= $ \rule{4cm}{0.4pt} \hfill
target $= $ \rule{4cm}{0.4pt}
\vspace{8pt}

\textbf{(b)} Now the same for a \emph{backward} branch: $PC = \code{0x30}$,
immediate $= \code{0xFFFC}$.

\vspace{4pt}
\code{NSA} $= $ \rule{4cm}{0.4pt} \hfill
offset $= $ \rule{4cm}{0.4pt} \hfill
target $= $ \rule{4cm}{0.4pt}
\vspace{8pt}

\textbf{(c)} Two students get part (a) wrong. One forgets the shift by two; the
other uses $PC$ instead of $PC+4$ as the base. What target does each of them
compute, and which of the two mistakes would be harder to notice? Why?

\answerbox{2.6cm}

\section*{Question 2 --- Interface and reuse}

\textbf{(a)} Write the VHDL \code{entity} for \code{NIAG\_Reg}, the program
counter. Include the \code{generic} for the reset vector.

\answerbox{4.0cm}

\textbf{(b)} Write the characteristic equation for $PC^{+}$:

\vspace{6pt}
$PC^{+} = $ \rule{10cm}{0.4pt}
\vspace{8pt}

Compare it with your Lab 2 register equation. In one sentence: what actually
changed?

\answerbox{1.4cm}

\section*{Question 3 --- Test plan}

\begin{center}
\begin{tabular}{@{}c@{\hskip 8pt}p{5.0cm}@{\hskip 8pt}p{7.6cm}@{}}
\toprule
\# & inputs & what it proves \\
\midrule
1 & \rule{4.8cm}{0.4pt} & \rule{7.4cm}{0.4pt} \\[8pt]
2 & \rule{4.8cm}{0.4pt} & \rule{7.4cm}{0.4pt} \\[8pt]
3 & \rule{4.8cm}{0.4pt} & \rule{7.4cm}{0.4pt} \\[8pt]
4 & \rule{4.8cm}{0.4pt} & \rule{7.4cm}{0.4pt} \\[8pt]
5 & \rule{4.8cm}{0.4pt} & \rule{7.4cm}{0.4pt} \\[8pt]
6 & \rule{4.8cm}{0.4pt} & \rule{7.4cm}{0.4pt} \\
\bottomrule
\end{tabular}
\end{center}

\textbf{Then answer both:}

\textbf{(a)} Suppose every branch offset in your test set is positive. Which
specific bug would you then be unable to detect?

\answerbox{1.6cm}

\textbf{(b)} \code{NSA} is combinational --- it changes as soon as \code{PC}
does, with no clock edge. Describe a vector that would prove this, and say what
would go wrong in the CPU if \code{NSA} were computed inside a clocked process
instead.

\answerbox{2.2cm}

\vfill
{\small\emph{Question 3 is the one that matters.} It is the same question as
Section 5 of your lab report.}

\end{document}
